\documentclass[10pt]{article}
\usepackage[margin=20mm]{geometry}
\usepackage{amsmath,amssymb,fontspec,unicode-math,xurl,hyperref}
\IfFontExistsTF{DejaVu Serif}{\setmainfont{DejaVu Serif}}{\setmainfont{Latin Modern Roman}}
\setmathfont{latinmodern-math.otf}
\setlength{\parindent}{0pt}\setlength{\parskip}{6pt}\sloppy
\begin{document}
Abstract: Three demonstrated mathematical problems affect the central theoretical presentation. The finite-volume domain assumption does not imply the Voronoi-weight or reward-continuity rates used in Lemma G.2; a positive-weight construction explicitly violates the claimed rate. Lemma G.4 invokes a curvature modulus for the old weights although its objective uses the new weights, and a bounded one-dimensional example disproves that invocation. Theorem 2.5 and the claimed uniform-weight equivalence also have inconsistent regularization normalization. Each concerns a core derivation or stated relationship and is assessed Major. Local corrections are established where possible. These results do not refute qualitative clone robustness on suitably specified domains, the Borda reversal example, or the paper’s general research motivation. Adjudication remains pending.\par
Author: AI-assisted technical review\par
Date: Initial review, 25 October 2026\par
Title: Technical Verification of
Clone-Robust AI Alignment\par
\# Identity, version, and inspected scope\par
The paper is Clone-Robust AI Alignment, by Ariel D. Procaccia, Benjamin Schiffer, and Shirley Zhang, ICML 2025, PMLR 267; identifier I25-43664, display 0295, OpenReview `vbw9hSlOaG` [Retained latex: \textbackslash{}cite\{paper\}]. The reviewed version is the staged 24-page proceedings PDF. Page references are PDF pages, agreeing with printed numbering.\par
The review examined the definitions and results in Sections 2–4, the main robustness proof in Appendix G, its convexity dependency in Appendix E, the normalization proof in Appendix F, the supporting impossibility and clone-reversal arguments, and the case-study design in Section 5 and Appendix I. The primary assistant independently derived the decisive examples below and inspected original source-page renderings. Auxiliary agents supplied mathematical challenges and checks of the supporting results. No attributable author repository was located in the completed source preparation; this is an availability limitation, not a claim that none exists. Embedded appendices were available, while external OpenReview attachment discovery was blocked.\par
\# Definitions and operative assumptions\par
The alternative space \(S\subset\mathbb R^d\) is assumed to have finite volume, and the observed set \(M\) is finite. Annotator rewards are \(K\)-Lipschitz in context distance; pairwise probabilities are population averages of Bradley–Terry–Luce probabilities. A representative dataset has exactly those pairwise probabilities. The clone-robustness definition asks for qualitative continuity when a new alternative approaches an existing one.\par
Equation reference: eq:weighted. Let \(\sigma(t)=(1+e^{-t})^{-1}\). Write the unweighted ordered-pair negative log likelihood as \[L_M(r)=-\sum_{x,y\in M}p_{xy}\log\sigma(r_x-r_y).\] Equation (1) minimizes \(F_\lambda(r)=L_M(r)+(\lambda/2)\sum_xr_x^2\). The weighted objective in Equation (2) is \[
 f_w(r)=-\sum_{x,y\in M}w_xw_yp_{xy}\log\sigma(r_x-r_y)
       +\frac{\lambda}{2}\sum_xw_xr_x^2,\] where \(w_x\) is the normalized volume of the Voronoi cell, splitting ties. The examples in Problems 1–3 all have positive cell weights, positive \(\lambda\), and valid representative probabilities. Thus none relies on a zero-weight ambiguity.\par
\# Substantive findings\par
Problem 1 (Finite volume does not give the rate asserted in Lemma G.2).  Assessment: Major; outcome: issue found.\par
\#\# Claim and dependency\par
Lemma G.2, p. 18, states \(|\widehat r_w^D(y)-r_1(y)|=O(\sqrt\epsilon)\) when a clone is within \(\epsilon\) of its original. Its Equation (13) first asserts that merged Voronoi weights change by \(O(\epsilon)\). Here \(r_1\) is the weighted optimizer after replacing the old weight of the cloned alternative by the combined weights of it and its clone, while retaining the old alternatives and probabilities. Theorem 4.2’s proof on pp. 16–18 uses this rate to choose \(\epsilon=c\delta^2\). The model on p. 3 assumes finite volume, not boundedness of \(S\).\par
\#\# Admissible domain and actual weight movement\par
Consider the fixed two-dimensional domain \[\begin{aligned}
 S&=([0,1]\times[-1,1])\cup T,\\
 T&=\left\{(u,v):v\ge e^2,\quad
       \frac34\le u\le\frac34+\frac1{v(\log v)^2}\right\}.\end{aligned}\] It is measurable with positive finite volume \(V=2+\int_{e^2}^{\infty}[v(\log v)^2]^{-1}dv=5/2\). Let \(M=\{a=(0,0),b=(1,0)\}\) and add \(a'=(0,\epsilon)\), where \(0<\epsilon<e^{-2}\). All three Voronoi cells have positive area within the rectangle. The distance between \(a'\) and \(a\) is exactly \(\epsilon\).\par
Equation reference: eq:q. Every tail point originally belongs to \(b\), because \(u\ge3/4>1/2\). The boundary between \(a'\) and \(b\) is \[u=\frac12+\epsilon v-\frac{\epsilon^2}{2}.\] For \(v\ge1/\epsilon\), this boundary is larger than 1, whereas the entire tail has \(u<1\). Thus all tail area above \(1/\epsilon\) moves into the combined \(a,a'\) cell. Adding \(a'\) cannot remove any old \(a\) cell from that combined cell. Consequently the increase \(q\) in the merged weight satisfies \[
 q\ge \frac1V\int_{1/\epsilon}^{\infty}\frac{dv}{v(\log v)^2}
   =\frac1{V\log(1/\epsilon)}.\] This is not \(O(\epsilon)\).\par
\#\# Reward calculation, not merely a geometric objection\par
Use one annotator with reward \(r^*(u,v)=(\log3)u\), globally Lipschitz with \(K=\log3\), and set \(\lambda=1\). Then \(p_{ab}=1/4\) and the clone has exactly the same preferences as \(a\). All comparison probabilities are \(1/4\), \(3/4\), or \(1/2\), so finite representative datasets exist.\par
For any two positive weights \(w_a+w_b=1\), the first-order conditions for [eq:weighted] (\#eq:weighted) give a fixed reward gap \(d=r_a-r_b\) satisfying \[d+\sigma(d)=\tfrac14,\qquad
 r_a=w_b d,\quad r_b=-w_a d.\] The unique root is \(d=-0.2001330666939489\ldots\ne0\). Increasing \(w_a\) by \(q\) therefore shifts both original rewards by \(-qd\). By [eq:q] (\#eq:q), \[|r_1(y)-\widehat r_w^D(y)|
 \ge\frac{|d|}{V\log(1/\epsilon)},\qquad y\in M.\] Dividing by \(\sqrt\epsilon\) diverges as \(\epsilon\downarrow0\). This directly falsifies Lemma G.2, with every relevant original assumption and positivity condition satisfied. The geometry and algebra are exact; `check.json` records supporting root and rate evaluations.\par
\#\# Counterevidence and bounded conclusion\par
The lower bound still tends to zero. The example refutes the asserted rate and the proof’s quadratic choice of neighborhood size, not the qualitative continuity in Definition 3.1. For finite-volume measurable domains, dominated convergence can support qualitative continuity of merged cells away from the null tie hyperplanes. A bounded-domain restriction would also change the geometry behind the claimed rate, but would be an additional assumption. No unconditional repaired version of Theorem 4.2 is certified here. The demonstrated false lemma is directly used in its central proof, supporting Major.\par
Problem 2 (Lemma G.4 uses the wrong weights in its curvature bound).  Assessment: Major; outcome: issue found.\par
\#\# Claim and complete bounded-domain check\par
On p. 21, Lemma G.4 invokes Lemma E.1 with \(\lambda\min_{z\in M}w_M(z)\) to convert an objective gap for \(f_2\) into a reward bound. But \(f_2\), defined on p. 17, uses the enlarged-set weights \(w_{M'}\). Those weights may approach zero even when the original minimum weight is fixed.\par
Take \(S=[0,1]\), \(M=\{0,1\}\), and clone \(x'=\epsilon=0.01\). The new weights in the order \((0,\epsilon,1)\) are \((0.005,0.5,0.495)\); the original minimum weight is \(0.5\). Let all annotator rewards be zero, all pair probabilities be \(1/2\), and \(\lambda=1\). Constant rewards satisfy Assumption 2.1 for any positive Lipschitz upper bound; finite representative datasets exist. The unique optimizer of \(f_2\) is \(r_2=0\).\par
For \(r=(1,0,0)\), direct evaluation of the ordered-pair objective gives, with \(w_0=0.005\), \[f_2(r)-f_2(0)
 =w_0(1-w_0)\log\cosh(1/2)+\frac{w_0}{2}
 =0.00309756967211743\ldots.\] Even the correctly halved strong-convexity lower bound using the old modulus would be \((1\cdot0.5/2)\|r-r_2\|_2^2=0.25\), which fails. The new-weight lower bound \(0.0025\) is compatible with the calculation. Hence the old modulus cannot justify the displayed implication. This witness has a bounded regular domain and strictly positive weights; it is separate from Problem 1.\par
\#\# Verified local replacement\par
Equation reference: eq:repair. There is a stronger local argument that avoids a minimum-weight bound. For two objectives sharing positive weights and \(\lambda\), divide their first-order conditions by each coordinate’s positive weight. For reward difference \(h=\widehat r-r_2\), subtraction gives \[\lambda h_i+\sum_j w_j\{\sigma(\widehat r_i-\widehat r_j)
                  -\sigma((r_2)_i-(r_2)_j)\}
       =\sum_jw_j(p'_{ij}-p^{(2)}_{ij}).\] At a coordinate maximizing \(h_i\), every sigmoid difference on the left is nonnegative. The minimum coordinate gives the reverse bound. Therefore \[
 \|\widehat r-r_2\|_\infty
 \le\frac{\max_{i,j}|p'_{ij}-p^{(2)}_{ij}|}{\lambda}
 \le\frac{K\epsilon}{4\lambda},\] using the sigmoid derivative bound \(1/4\) and the original Lipschitz assumption. This holds uniformly despite shrinking positive clone weights and was independently challenged.\par
The original invocation is invalid in a step directly proving Theorem 4.2, hence Major. Equation [eq:repair] (\#eq:repair) locally repairs this part for positive weights; it does not repair Problem 1’s asserted rate or settle all domain issues.\par
Problem 3 (The moment identity and uniform-weight equivalence use inconsistent normalization).  Assessment: Major; outcome: issue found.\par
\#\# Original statements and independent derivation\par
Theorem 2.5, p. 4, claims \(\operatorname{AWR}_D(x)=\lambda\widehat r_x+\widehat{\operatorname{AWR}}(x)\) for Equation (1), with average win rate defined using \(1/m\). The paragraph on p. 6 also claims that uniform weights make Equation (2) a scaled version of Equation (1), with the same optimizer.\par
For Equation (1), differentiation instead gives \[0=\lambda r_x+\sum_y\{\sigma(r_x-r_y)-p_{xy}\},\] and hence \[\operatorname{AWR}_D(x)=\frac{\lambda}{m}r_x+
                  \widehat{\operatorname{AWR}}(x).\] The one-sentence proof in Appendix F substitutes weights 1 into the weighted identity, but does not divide the regularizer term when converting sums into averages. With uniform \(w_x=1/m\), Equation (2) is \[f_w(r)=\frac1{m^2}L_M(r)+\frac{\lambda}{2m}\sum_xr_x^2
       =\frac1{m^2}F_{m\lambda}(r),\] so it matches the standard optimizer at regularization \(m\lambda\), not \(\lambda\).\par
\#\# Explicit legitimate discrepancy\par
Let \(S=[0,1]\), \(M=\{0,1\}\), \(p_{ab}=3/4\), and \(\lambda=1\). The Voronoi weights are both \(1/2\). A single Lipschitz annotator with reward gap \(\log3\) generates the representative probabilities. Strict convexity and symmetry imply optimizers \((a,-a)\) and \((b,-b)\), where \[\sigma(2a)+a=\tfrac34,\qquad
 \sigma(2b)+2b=\tfrac34.\] Independent bisection gives \(a=0.1671800993781829\) and \(b=0.10006653334697452\). Theorem 2.5’s residual for alternative \(a\) is \(-0.08359004968909145\); the corrected \(\lambda/m\) identity has residual zero to the recorded precision. The analytic stationarity equations establish the error independently of the numerical solver.\par
These false statements concern the claimed preservation of the baseline’s interpretation and its equality to the new method in uniform instances. This supports Major for the affected mathematical relationship, without claiming empirical score loss. Retuning \(\lambda_{\rm standard}=m\lambda_{\rm weighted}\) restores the comparison. Alternatively, replacing Equation (1)’s loss by \(L_M/m\) would make the printed moment coefficient and same-\(\lambda\) uniform comparison consistent, but changes the defined baseline. Only this algebraic correction is verified; historical experiments were not rerun.\par
The ranking relationship survives a positive coefficient \(\lambda/m\), so this problem does not invalidate the Borda-based clone reversal in Theorem 3.2. The weighted identity of Theorem 4.4 itself is valid for positive weights.\par
\# Surviving results, domain qualifications, and empirical limits\par
The two-population indistinguishability argument in Theorem 2.3 was checked within its variable reward-scale setting. The Borda example in Appendix D was independently recalculated: the original scores are approximately \((1.589919,1.5,1.410081)\) and the cloned scores \((2.040504,2.139334,1.910081,1.910081)\), retaining the winner reversal from \(a\) to \(b\). Theorem 4.6’s finite-sum/integral substitution was checked under measurable positive finite-volume conditions.\par
The domain also needs a rule for zero-volume cells. For example, \(S=[0,1]\cup\{2\}\) and \(M=\{0,1,2\}\) satisfy finite positive volume but assign weight zero to 2. The objective then leaves its reward unconstrained. An accumulating measure-zero sequence of potential clones has the same issue. This establishes nonuniqueness under the bare domain assumptions, not failure of every possible tie-breaking rule. Positive-cell assumptions or a specified extension are required before the proposed local repairs become a fully defined global algorithm. The intended restriction or rule remains unresolved.\par
The case study is explicitly synthetic and departs from complete representative comparisons. Appendix I describes 1000 comparisons per dataset, separately fitted 1500-component PCA, a small neural network, 20 training runs, and Monte Carlo weight estimation. No released dataset, implementation, regularization value, seeds, or historical model configuration was available for verification. Response-level standard-deviation bars and dataset size alone do not establish that a shift cannot be due to generation/training variability. Run-level uncertainty and the relationship between separate PCA coordinates and one fixed \(S\) remain Needs Evidence. No paid API calls, model training, or empirical reproduction was attempted.\par
\# Overall impact and state\par
The overall assessment is Major. Two independent flaws affect the proof of the central robustness theorem, and a further normalization error changes a claimed exact relationship between the proposed and baseline objectives. The report distinguishes these from a refutation of qualitative clone robustness: the rate counterexample still converges, and a valid local replacement is given for G.4. A complete repaired theorem requires explicit domain conditions and a revised continuity argument.\par
Item | Assessment | Repair verification | Adoption / remaining work\par
G.2 continuity rate | Major | Rate refuted; no full repair | Unknown; revise domain/proof\par
G.4 curvature step | Major | Positive-weight local bound proved | Unknown; global domain needed\par
Normalization | Major | Correct identities verified | Unknown; reconcile baseline\par
Case study/domain rule | Needs Evidence | Scope documented | Data and intended rule needed\par
Scientific adjudication is pending and `human\_reviews` is empty. No author adoption or human endorsement is claimed.\par
9 A. D. Procaccia, B. Schiffer, and S. Zhang. Clone-Robust AI Alignment. ICML, PMLR 267 (2025). https://proceedings.mlr.press/v267/procaccia25a.html. Staged 24-page proceedings PDF, OpenReview `vbw9hSlOaG`; embedded Appendices A–I.\par
\# Evidence and reproducibility\par
Original PDF SHA-256: . Manifest SHA-256: . All four staged inputs are preserved; no author-code archive is available.\par
Reviewer-written `check.py` records 180-step bisection, direct log-cosh evaluation, and domain/rate diagnostics. Analytic arguments establish the claims independently of those numerical checks. Original pp. 4, 6, 18, and 21 were rendered and visually inspected.\par
Native auxiliaries supplied supporting-proof/source checks (`a1.json`) and mathematical challenge (`a2.json`). The primary assistant independently constructed the continuous-tail example, checked calculations, and owns integration, compilation, and all-page visual QA. Unexposed model labels and chat URL remain null; no human review is claimed.\par
The frozen template is compiled with available isolated `pdflatex`, shell escape disabled. Actual commands, logs, final TeX/PDF hashes, coverage, and visual records are retained. No author program, paid API, installation, additional download, author contact, or publication was used.\par
\end{document}
